\documentclass[12pt, a4paper]{article}

% --- Pakete ---
\usepackage[ngerman]{babel}
\usepackage[utf8]{inputenc}
\usepackage[T1]{fontenc}
\usepackage{lmodern}
\usepackage{geometry}
\usepackage{graphicx}
\usepackage{hyperref}
\usepackage{listings}
\usepackage[table]{xcolor}
\usepackage{booktabs}
\usepackage{enumitem}
\usepackage{parskip}
\usepackage{fancyhdr}
\usepackage{tcolorbox}
\usepackage{tabularx}
\usepackage{tikz}

\tcbuselibrary{skins, breakable}

\geometry{left=2.5cm, right=2.5cm, top=2.8cm, bottom=2.5cm, headheight=25pt}

% --- Farben ---
\definecolor{primary}{RGB}{41, 98, 155}
\definecolor{accent}{RGB}{231, 76, 60}
\definecolor{success}{RGB}{39, 174, 96}
\definecolor{infoblue}{RGB}{52, 152, 219}
\definecolor{warnorg}{RGB}{243, 156, 18}
\definecolor{codebg}{RGB}{248, 249, 250}
\definecolor{codeborder}{RGB}{222, 226, 230}
\definecolor{codegray}{RGB}{108, 117, 125}
\definecolor{sectioncolor}{RGB}{41, 98, 155}

% --- Abschnitt-Stil ---
\usepackage{titlesec}
\titleformat{\section}
  {\normalfont\Large\bfseries\color{sectioncolor}}
  {\thesection}{1em}{}[\color{sectioncolor}\titlerule]
\titleformat{\subsection}
  {\normalfont\large\bfseries\color{sectioncolor!80!black}}
  {\thesubsection}{0.8em}{}
\titleformat{\subsubsection}
  {\normalfont\normalsize\bfseries\color{sectioncolor!60!black}}
  {\thesubsubsection}{0.6em}{}

% --- tcolorbox Stile ---
\newtcolorbox{infobox}[1][]{
    colback=infoblue!5,
    colframe=infoblue!70,
    fonttitle=\bfseries,
    coltitle=white,
    colbacktitle=infoblue,
    boxrule=0.5pt,
    arc=3pt,
    left=8pt, right=8pt, top=6pt, bottom=6pt,
    title={#1},
    breakable,
}

\newtcolorbox{warnbox}[1][]{
    colback=warnorg!5,
    colframe=warnorg!70,
    fonttitle=\bfseries,
    coltitle=white,
    colbacktitle=warnorg,
    boxrule=0.5pt,
    arc=3pt,
    left=8pt, right=8pt, top=6pt, bottom=6pt,
    title={#1},
    breakable,
}

\newtcolorbox{successbox}[1][]{
    colback=success!5,
    colframe=success!70,
    fonttitle=\bfseries,
    coltitle=white,
    colbacktitle=success,
    boxrule=0.5pt,
    arc=3pt,
    left=8pt, right=8pt, top=6pt, bottom=6pt,
    title={#1},
    breakable,
}

\newtcolorbox{stepbox}[1][]{
    colback=primary!5,
    colframe=primary!70,
    fonttitle=\bfseries,
    coltitle=white,
    colbacktitle=primary,
    boxrule=0.5pt,
    arc=3pt,
    left=8pt, right=8pt, top=6pt, bottom=6pt,
    title={#1},
    breakable,
}

% --- Header/Footer ---
\pagestyle{fancy}
\fancyhf{}
\fancyhead[L]{\color{sectioncolor}\textbf{KuchenApp} \color{codegray}-- Benutzerdokumentation}
\fancyhead[R]{\color{codegray}\thepage}
\renewcommand{\headrulewidth}{0.4pt}
\renewcommand{\headrule}{\hbox to\headwidth{\color{sectioncolor}\leaders\hrule height \headrulewidth\hfill}}

% --- Hyperlinks ---
\hypersetup{
    colorlinks=true,
    linkcolor=primary,
    urlcolor=infoblue,
}

% =====================================================================
\begin{document}

% --- Titelseite ---
\begin{titlepage}
\centering
\begin{tikzpicture}[remember picture, overlay]
    \fill[primary] (current page.north west) rectangle ([yshift=-6cm]current page.north east);
    \fill[primary!80!black] ([yshift=-6cm]current page.north west) rectangle ([yshift=-6.3cm]current page.north east);
\end{tikzpicture}

\vspace*{1cm}
{\color{white}\Huge\textbf{KuchenApp}}\\[0.4cm]
{\color{white!85}\LARGE Benutzerdokumentation}\\[4.5cm]

{\Large\color{primary} Anleitung zur Bedienung der\\[0.3cm]
KuchenApp fuer Endanwender}\\[2.5cm]

\begin{tcolorbox}[
    colback=codebg, colframe=codeborder,
    width=10cm, arc=4pt, boxrule=0.5pt,
    halign=center,
]
\begin{tabular}{rl}
    \textbf{Anwendung:} & KuchenApp \\
    \textbf{Plattformen:} & Linux \& Windows \\
    \textbf{Zielgruppe:} & Klassensprecher / Lehrkraft \\
    \textbf{Stand:} & \today \\
\end{tabular}
\end{tcolorbox}

\vfill
{\color{codegray}\small Berufsschule -- Fachinformatiker Anwendungsentwicklung}
\end{titlepage}

\tableofcontents
\newpage

% =====================================================================
\section{Einleitung}

Diese Benutzerdokumentation richtet sich an die Anwender der \textbf{KuchenApp}. Sie
beschreibt Schritt fuer Schritt, wie die Anwendung installiert, gestartet und
bedient wird. Technische Hintergrund\-informationen sind in der separaten
\emph{Projektdokumentation} zu finden.

\begin{infobox}[Was kann die KuchenApp?]
Die KuchenApp hilft dir dabei, eine Liste an Personen zu verwalten, die fuer einen
bestimmten Termin Kuchen mitbringen. Mit einem Klick erstellt sie dir eine
Erinnerungs-E-Mail und verschickt sie ueber dein IServ-Konto an alle Empfaenger.
\end{infobox}

% =====================================================================
\section{Installation}

\subsection{Systemvoraussetzungen}

\begin{itemize}[leftmargin=1.5em]
    \item Betriebssystem: Windows 10/11 oder Linux (64-Bit)
    \item Internetverbindung (nur fuer den Mail-Versand erforderlich)
    \item Gueltiges IServ-Konto der Schule
\end{itemize}

\subsection{Installation unter Windows}

\begin{stepbox}[Schritt fuer Schritt]
\begin{enumerate}[leftmargin=1.5em]
    \item Lade die aktuelle Version von der GitHub-Releases-Seite herunter:\\
          \url{https://github.com/Albazos/Kuchen/releases}
    \item Entpacke die ZIP-Datei in einen Ordner deiner Wahl.
    \item Starte die Anwendung per Doppelklick auf \texttt{KuchenApp.exe}.
\end{enumerate}
\end{stepbox}

\subsection{Installation unter Linux}

\begin{stepbox}[Schritt fuer Schritt]
\begin{enumerate}[leftmargin=1.5em]
    \item Lade die Linux-Version von der Releases-Seite herunter.
    \item Entpacke die ZIP-Datei in einen Ordner deiner Wahl.
    \item Mache die Datei ausfuehrbar: \texttt{chmod +x KuchenApp}
    \item Starte die Anwendung mit: \texttt{./KuchenApp}
\end{enumerate}
\end{stepbox}

\begin{warnbox}[Hinweis zur Windows-Sicherheitswarnung]
Beim ersten Start erscheint unter Windows ggf. der \emph{SmartScreen}-Warnhinweis,
da die Anwendung nicht signiert ist. Klicke auf \emph{Weitere Informationen}
$\rightarrow$ \emph{Trotzdem ausfuehren}, um die App zu starten.
\end{warnbox}

% =====================================================================
\section{Das Hauptfenster im Ueberblick}

Nach dem Start zeigt die KuchenApp das Hauptfenster mit folgenden Bereichen:

\begin{center}
\begin{tabularx}{\linewidth}{l X}
\toprule
\textbf{Bereich} & \textbf{Funktion} \\
\midrule
\textbf{Tabelle} & Zeigt die aktuelle Kuchen-Liste an. Zellen koennen direkt bearbeitet werden. \\
\textbf{Suche} & Filtert die Tabelle nach dem eingegebenen Begriff. \\
\textbf{Sort} & Auswahl der Spalte, nach der sortiert wird (auf-/absteigend). \\
\textbf{Buttons} & Aktionen wie Importieren, Speichern, Hinzufuegen, Loeschen, Mailliste, Senden. \\
\textbf{Hinweisdialoge} & Zeigen Erfolgs-, Warn- und Fehlermeldungen als Dialogfenster an. \\
\bottomrule
\end{tabularx}
\end{center}

% =====================================================================
\section{Arbeiten mit der Kuchen-Liste}

\begin{warnbox}[Hinweis zu den mitgelieferten Daten]
Im Repository sollen nur neutrale Template- und Standarddateien versioniert
werden. Persoenliche oder im Schulalltag verwendete Daten sollten lokal liegen
und nicht mit Git veroeffentlicht werden.
\end{warnbox}

Beim ersten Start werden die mitgelieferten Standarddateien geladen. Sie enthalten
nur die benoetigten Kopfzeilen. Eigene Klassen- und Kuchendaten koennen entweder
direkt in der App eingetragen oder ueber CSV-Dateien importiert werden.

\subsection{CSV-Datei importieren}

\begin{stepbox}[Datei oeffnen]
\begin{enumerate}[leftmargin=1.5em]
    \item Klicke auf den Button \textbf{Import}.
    \item Waehle im Dateidialog die gewuenschte CSV-Datei aus.
    \item Bestaetige mit \emph{Oeffnen}. Die Tabelle wird mit den Daten gefuellt.
\end{enumerate}
\end{stepbox}

\begin{infobox}[CSV-Format]
Eine gueltige CSV-Datei besteht aus einer Kopfzeile und beliebig vielen
Datenzeilen, jeweils kommagetrennt. Beispiel:
\begin{center}
\texttt{Name,Datum,Status} \\
\texttt{Max Mustermann,01.05.2026,offen} \\
\texttt{Erika Beispiel,08.05.2026,zugesagt}
\end{center}
\end{infobox}

\subsection{Eintraege bearbeiten, hinzufuegen oder loeschen}

\begin{itemize}[leftmargin=1.5em]
    \item \textbf{Bearbeiten:} Klicke doppelt in eine Zelle und gib den neuen Wert ein.
    \item \textbf{Hinzufuegen:} Klicke auf \textbf{Add}, um eine leere Zeile am Ende der Tabelle einzufuegen.
    \item \textbf{Loeschen:} Markiere eine oder mehrere Zeilen und klicke auf \textbf{Delete}.
\end{itemize}

\subsection{Sortieren und Suchen}

\begin{itemize}[leftmargin=1.5em]
    \item Im Bereich \textbf{Sort} eine Spalte sowie \emph{Ascending} oder
          \emph{Descending} auswaehlen -- die Tabelle wird sofort umsortiert.
    \item Im Suchfeld einen Begriff eingeben -- die Tabelle zeigt nur noch
          passende Zeilen an. Ist eine Sortier-Spalte aktiv, wird nur in dieser
          Spalte gesucht; andernfalls wird die gesamte Tabelle durchsucht.
\end{itemize}

\subsection{Aenderungen speichern}

Klicke auf \textbf{Save}, um die aktuellen Daten zurueck in die CSV-Datei zu
schreiben. Beim Beenden der Anwendung wird ausserdem gefragt, ob ungespeicherte
Aenderungen gesichert werden sollen.

\begin{warnbox}[Datenverlust vermeiden]
Speichere regelmaessig zwischen. Beendest du die App ueber das X des Fensters
ohne Rueckfrage des Betriebssystems zu beachten, gehen ungesicherte Aenderungen
verloren.
\end{warnbox}

% =====================================================================
\section{Mailliste verwalten}

\subsection{Empfaenger hinzufuegen, aendern oder entfernen}

\begin{stepbox}[Mailliste oeffnen]
\begin{enumerate}[leftmargin=1.5em]
    \item Klicke im Hauptfenster auf \textbf{Open Mail List}.
    \item Es oeffnet sich ein separater Dialog mit allen aktuellen Empfaengern.
    \item Pflege die Eintraege wie in der Hauptliste (bearbeiten, hinzufuegen,
          loeschen) und schliesse den Dialog ueber den Schliessen-Button.
\end{enumerate}
\end{stepbox}

\begin{infobox}[Empfaengerformat]
Jede Zeile enthaelt einen Namen und eine gueltige IServ-E-Mail-Adresse, z.\,B.
\texttt{vorname.nachname@meine-schule.de}.
\end{infobox}

% =====================================================================
\section{Erinnerungs-Mails versenden}

\subsection{Versand starten}

\begin{stepbox}[Mails senden]
\begin{enumerate}[leftmargin=1.5em]
    \item Stelle sicher, dass die Kuchen-Liste und die Mailliste aktuell sind.
    \item Klicke auf \textbf{Send Mail}.
    \item Bestaetige die Sicherheitsabfrage \emph{"`Mail an X Empfaenger
          senden?"'} mit \textbf{Ja}.
    \item Im folgenden Login-Dialog deinen IServ-Benutzernamen und das
          dazugehoerige Passwort eingeben.
    \item Auf \textbf{Login / Senden} klicken -- die Mails werden ueber den
          IServ-Mailserver verschickt.
\end{enumerate}
\end{stepbox}

\begin{successbox}[Erfolgreicher Versand]
Bei erfolgreichem Versand erscheint ein Informationsdialog mit der Meldung
\emph{"`Mails sent successfully."'}.
\end{successbox}

\subsection{Was wird verschickt?}

Die App generiert automatisch eine HTML-E-Mail mit der aktuellen Kuchen-Liste
als formatierte Tabelle. Jeder Empfaenger erhaelt dieselbe Uebersicht und kann
so seinen Eintrag direkt erkennen.

\begin{warnbox}[Hinweis zum Passwort]
Dein IServ-Passwort wird ausschliesslich fuer den aktuellen Versand verwendet
und \textbf{nicht} dauerhaft gespeichert. Du musst es bei jedem Versand neu
eingeben.
\end{warnbox}

% =====================================================================
\section{Haeufige Fragen (FAQ)}

\subsection{Die Anwendung startet nicht}
Pruefe, ob du das Archiv vollstaendig entpackt hast und die Anwendung aus dem
entpackten Ordner heraus startest. Unter Linux muss die Datei ausfuehrbar sein
(\texttt{chmod +x KuchenApp}).

\subsection{Beim Import passiert nichts}
Stelle sicher, dass die Datei tatsaechlich im CSV-Format vorliegt und eine
Kopfzeile besitzt. Im Fehlerfall erscheint ein Dialog mit einer Meldung wie
\emph{"`Import failed."'} oder \emph{"`No file selected."'}.

\subsection{Die Mails werden nicht verschickt}
Pruefe folgende Punkte:
\begin{itemize}[leftmargin=1.5em]
    \item Besteht eine Internetverbindung?
    \item Sind Benutzername und Passwort korrekt?
    \item Enthaelt die Mailliste mindestens einen gueltigen Empfaenger?
\end{itemize}
Erscheint die Meldung \emph{"`Failed to send mails."'}, wiederhole den Login.

\subsection{Wie kann ich Aenderungen rueckgaengig machen?}
Die App besitzt keine Undo-Funktion. Solange du die Datei \emph{nicht}
gespeichert hast, kannst du den Import ueber \textbf{Import} erneut ausfuehren,
um die Originaldaten wiederherzustellen.

% =====================================================================
\section{Support}

Bei Fehlern oder Verbesserungsvorschlaegen erstelle bitte ein Issue im
GitHub-Repository:\\
\url{https://github.com/Albazos/Kuchen/issues}

\end{document}
